% ch5-eigenvalues.tex — Chapter 5: Eigenvalues, Eigenvectors, and Diagonalization
% Definitions, characteristic polynomial, Cayley–Hamilton, diagonalizability,
% minimal polynomial, applications (A^n, Fibonacci, Markov chains, ODEs).

\chapter{Eigenvalues, Eigenvectors, and Diagonalization}\label{ch:eigenvalues}

% ============================================================
% Motivating introduction
% ============================================================

Among all the questions one can ask about a linear map, perhaps the
most fundamental is this: \emph{does there exist a basis in which the
map takes a particularly simple form?}  The simplest form one can hope
for is a diagonal matrix, and the search for such a basis leads
directly to the theory of eigenvalues and eigenvectors.

This theory is not merely an exercise in abstraction.  Eigenvalues
appear throughout mathematics and its applications:
\begin{itemize}
  \item \textbf{Differential equations.}  The solutions of a linear
    system $\vec{x}'(t) = A\vec{x}(t)$ are determined by the
    eigenvalues of~$A$; exponential growth, decay, and oscillation all
    correspond to different eigenvalue configurations.
  \item \textbf{Markov chains.}  The long-run behaviour of a stochastic
    process is governed by the eigenvalue~$1$ and its eigenvector (the
    steady-state distribution).
  \item \textbf{Vibrations and resonance.}  The natural frequencies of a
    vibrating string, membrane, or mechanical system are the square
    roots of the eigenvalues of the associated stiffness matrix.
  \item \textbf{Principal component analysis (PCA).}  In statistics and
    data science, the principal components are the eigenvectors of the
    covariance matrix, and the eigenvalues measure the variance along
    each component.
\end{itemize}

\begin{center}
\begin{tikzpicture}[>=Stealth, scale=1.4]
  % --- Eigenvector illustration: stretching and flipping ---
  % Original basis vectors
  \fill[blue!10] (-0.3,-0.3) rectangle (2.5,2.5);
  \draw[->, thick, blue!70!black] (0,0) -- (2,1)
    node[right] {$\vec{v}$};
  \draw[->, thick, red!70!black] (0,0) -- (0.5,2)
    node[above] {$\vec{w}$};

  % Arrow for the map
  \draw[->, thick, dashed, gray] (3,1) to[bend left=15]
    node[midway,above] {\small $f$} (5,1);

  % Images
  \begin{scope}[shift={(6,0)}]
    \fill[blue!10] (-1.3,-0.8) rectangle (4.3,2.5);
    % Eigenvector stretched by 2
    \draw[->, thick, blue!70!black] (0,0) -- (4,2)
      node[right] {$f(\vec{v}) = 2\vec{v}$};
    % Eigenvector flipped (eigenvalue -1)
    \draw[->, thick, red!70!black] (0,0) -- (-0.5,-2)
      node[below] {$f(\vec{w}) = -\vec{w}$};
    \draw[->, thick, red!70!black, dashed, opacity=0.35] (0,0) -- (0.5,2)
      node[above, opacity=0.5] {\small (original $\vec{w}$)};
  \end{scope}

  % Labels
  \node[below] at (1.2,-0.5) {\small Eigenvector $\vec{v}$: eigenvalue $\lambda=2$};
  \node[below] at (7.5,-0.5) {\small Eigenvector $\vec{w}$: eigenvalue $\lambda=-1$};
\end{tikzpicture}
\captionof{figure}{An endomorphism acts on eigenvectors by scaling:
  $\vec{v}$ is stretched by factor~$2$, and $\vec{w}$ is flipped
  (eigenvalue~$-1$).}
\label{fig:eigenvector-intro}
\end{center}

Throughout this chapter, $\K$ denotes a field (typically $\R$ or~$\C$),
$E$ denotes a finite-dimensional $\K$-vector space with $\dim E = n$,
and $f\in\End(E)$ is an endomorphism of~$E$.  We shall use results
from 
ef{ch:linear-maps} freely, especially the notions of kernel,
image, determinant, and matrix representation.


% ============================================================
\section{Eigenvalues and eigenvectors}\label{sec:eigenvalues-def}
% ============================================================

\begin{definition}{Eigenvalue and eigenvector of an endomorphism}%
{def:eigenvalue-endomorphism}
\index{eigenvalue}\index{eigenvector}
Let $f\in\End(E)$.  A scalar $\lambda\in\K$ is called an
\emph{eigenvalue} of~$f$ if there exists a nonzero vector $v\in E$
such that
\[
  f(v) = \lambda v.
\]
Any such nonzero vector~$v$ is called an \emph{eigenvector} of~$f$
associated with (or belonging to) the eigenvalue~$\lambda$.
\end{definition}

\begin{remark}{The zero vector is never an eigenvector}{rem:zero-not-eigen}
By convention, the zero vector is \emph{not} an eigenvector, even
though $f(\vzero) = \lambda\vzero$ holds for every~$\lambda$.  This
convention ensures that the set of eigenvectors associated with a
given eigenvalue, together with~$\vzero$, forms a subspace.
\end{remark}

\begin{definition}{Eigenvalue and eigenvector of a matrix}%
{def:eigenvalue-matrix}
\index{eigenvalue!of a matrix}\index{eigenvector!of a matrix}
Let $A\in\Mat_n(\K)$.  A scalar $\lambda\in\K$ is an eigenvalue
of~$A$ if there exists a nonzero column vector
$X\in\K^n$ such that $AX = \lambda X$.  Such an~$X$ is an
eigenvector of~$A$ for~$\lambda$.
\end{definition}

\begin{proposition}{Reformulation via the kernel}%
{prop:eigenvalue-kernel}
Let $f\in\End(E)$ and $\lambda\in\K$.  Then $\lambda$ is an
eigenvalue of~$f$ if and only if
\[
  \Ker(f - \lambda\Id) \neq \set{\vzero}.
\]
Equivalently, for $A\in\Mat_n(\K)$, $\lambda$ is an eigenvalue of~$A$
if and only if $\det(A - \lambda I_n) = 0$.
\end{proposition}

\begin{proof}
We have $f(v) = \lambda v$ if and only if $(f - \lambda\Id)(v) = \vzero$,
i.e.\ $v\in\Ker(f - \lambda\Id)$.  A nonzero such~$v$ exists precisely
when $\Ker(f - \lambda\Id) \neq \set{\vzero}$, which is equivalent to
$f - \lambda\Id$ being non-injective, hence non-invertible (in finite
dimension).  For a matrix~$A$ this means $\det(A - \lambda I_n) = 0$.
\end{proof}

\begin{example}{Eigenvalues of a $2\times 2$ matrix}{ex:eigen-2x2}
Let $A = \begin{pmatrix} 3 & 1 \\ 0 & 2 \end{pmatrix}$.  We solve
$\det(A - \lambda I_2) = 0$:
\[
  \det\begin{pmatrix} 3-\lambda & 1 \\ 0 & 2-\lambda \end{pmatrix}
  = (3-\lambda)(2-\lambda) = 0,
\]
so $\lambda_1 = 3$ and $\lambda_2 = 2$.

For $\lambda_1 = 3$: $(A - 3I_2)X = \vzero$ gives
$\begin{pmatrix} 0 & 1 \\ 0 & -1 \end{pmatrix}
\begin{pmatrix} x \\ y \end{pmatrix} = \vzero$,
hence $y = 0$ and $x$ is free.  An eigenvector is
$\begin{psmallmatrix} 1 \\ 0 \end{psmallmatrix}$.

For $\lambda_2 = 2$: $(A - 2I_2)X = \vzero$ gives
$\begin{pmatrix} 1 & 1 \\ 0 & 0 \end{pmatrix}
\begin{pmatrix} x \\ y \end{pmatrix} = \vzero$,
hence $x = -y$.  An eigenvector is
$\begin{psmallmatrix} -1 \\ 1 \end{psmallmatrix}$.
\end{example}

\begin{proposition}{Eigenvectors for distinct eigenvalues are linearly
independent}{prop:distinct-eigenvalues-li}
\index{eigenvector!linear independence}
Let $f\in\End(E)$ and let $\lambda_1,\ldots,\lambda_r$ be pairwise
distinct eigenvalues of~$f$, with associated eigenvectors
$v_1,\ldots,v_r$.  Then $(v_1,\ldots,v_r)$ is linearly independent.
\end{proposition}

\begin{proof}
By induction on~$r$.  The case $r=1$ is clear since $v_1\neq\vzero$.

Suppose the result holds for $r-1$ eigenvalues and consider a
relation $\sum_{i=1}^{r}\alpha_i v_i = \vzero$.  Apply $f$:
\[
  \sum_{i=1}^{r}\alpha_i \lambda_i v_i = \vzero.
\]
Subtract $\lambda_r$ times the original relation:
\[
  \sum_{i=1}^{r-1}\alpha_i(\lambda_i - \lambda_r) v_i = \vzero.
\]
By the induction hypothesis, $\alpha_i(\lambda_i - \lambda_r) = 0$
for all $i = 1,\ldots,r-1$.  Since the eigenvalues are distinct,
$\lambda_i - \lambda_r \neq 0$, so $\alpha_i = 0$ for
$i = 1,\ldots,r-1$.  Substituting back gives $\alpha_r v_r = \vzero$,
whence $\alpha_r = 0$.
\end{proof}

\begin{corollary}{Bound on the number of eigenvalues}%
{cor:eigenvalue-bound}
An endomorphism of an $n$-dimensional space has at most $n$ distinct
eigenvalues.
\end{corollary}

\begin{proof}
If there were $n+1$ distinct eigenvalues, the corresponding
eigenvectors would form a linearly independent family of $n+1$ vectors
in an $n$-dimensional space, contradicting the definition of dimension.
\end{proof}


% ============================================================
\section{Eigenspaces}\label{sec:eigenspaces}
% ============================================================

\begin{definition}{Eigenspace}{def:eigenspace}
\index{eigenspace}
Let $f\in\End(E)$ and $\lambda\in\K$ be an eigenvalue of~$f$.  The
\emph{eigenspace} of~$f$ associated with~$\lambda$ is
\[
  E_\lambda(f) \deff \Ker(f - \lambda\Id)
  = \setcond{v\in E}{f(v) = \lambda v}.
\]
For a matrix $A\in\Mat_n(\K)$, we write $E_\lambda(A) = \Ker(A-\lambda I_n)$.
\end{definition}

\begin{proposition}{Eigenspaces are subspaces}{prop:eigenspace-subspace}
\index{eigenspace!is a subspace}
For every eigenvalue~$\lambda$ of~$f$, the eigenspace $E_\lambda(f)$
is a vector subspace of~$E$ of dimension at least~$1$.
\end{proposition}

\begin{proof}
Since $f - \lambda\Id$ is a linear map, its kernel $\Ker(f - \lambda\Id)$
is a subspace of~$E$ (the kernel of any linear map is a subspace).
It contains at least one nonzero vector (the eigenvector guaranteed by
the definition of eigenvalue), so $\dim E_\lambda(f) \geq 1$.
\end{proof}

\begin{proposition}{Sum of eigenspaces is direct}{prop:eigenspace-direct-sum}
\index{eigenspace!direct sum}
Let $\lambda_1,\ldots,\lambda_r$ be pairwise distinct eigenvalues
of~$f$.  Then the sum $E_{\lambda_1}(f) + \cdots + E_{\lambda_r}(f)$
is a direct sum:
\[
  E_{\lambda_1}(f) + \cdots + E_{\lambda_r}(f)
  = E_{\lambda_1}(f) \oplus \cdots \oplus E_{\lambda_r}(f).
\]
\end{proposition}

\begin{proof}
Suppose $v_1 + \cdots + v_r = \vzero$ with $v_i\in E_{\lambda_i}(f)$.
We must show each $v_i = \vzero$.  If some $v_i \neq \vzero$, let
$S = \setcond{i}{v_i \neq \vzero}$.  Then the nonzero vectors
$(v_i)_{i\in S}$ are eigenvectors for pairwise distinct eigenvalues,
hence linearly independent by

ef{prop:prop:distinct-eigenvalues-li}.  But
$\sum_{i\in S} v_i = \vzero$ (the terms with $v_i = \vzero$ contribute
nothing), contradicting linear independence.  Therefore $S = \emptyset$
and all $v_i = \vzero$.
\end{proof}


% ============================================================
\section{Characteristic polynomial}\label{sec:char-poly}
% ============================================================

\begin{definition}{Characteristic polynomial of a matrix}%
{def:char-poly-matrix}
\index{characteristic polynomial!of a matrix}
Let $A\in\Mat_n(\K)$.  The \emph{characteristic polynomial} of~$A$ is
\[
  \chi_A(\lambda) \deff \det(A - \lambda I_n) \in \K[\lambda].
\]
It is a polynomial of degree~$n$ in~$\lambda$.
\end{definition}

\begin{remark}{Sign convention}{rem:sign-convention}
Some authors define $\chi_A(\lambda) = \det(\lambda I_n - A)$
instead, which differs from our convention by a factor of~$(-1)^n$.
With our convention, the leading term is $(-1)^n\lambda^n$; with the
other convention it is $\lambda^n$.  Both conventions are standard.
The eigenvalues (roots of~$\chi_A$) are the same either way.
\end{remark}

\begin{proposition}{Properties of the characteristic polynomial}%
{prop:char-poly-properties}
Let $A\in\Mat_n(\K)$ with characteristic polynomial
$\chi_A(\lambda) = (-1)^n\lambda^n + c_{n-1}\lambda^{n-1} + \cdots + c_0$.
Then:
\begin{enumerate}[label=(\roman*)]
  \item $\deg\chi_A = n$ and the leading coefficient is $(-1)^n$.
  \item $c_{n-1} = (-1)^{n-1}\tr(A)$.
  \item $c_0 = \chi_A(0) = \det(A)$.
  \item The eigenvalues of~$A$ are exactly the roots of~$\chi_A$
    in~$\K$.
  \item Similar matrices have the same characteristic polynomial:
    if $B = \inv{P}AP$ then $\chi_B = \chi_A$.
\end{enumerate}
\end{proposition}

\begin{proof}
(i) The determinant of $A - \lambda I_n$ is a sum over permutations
$\sigma\in S_n$ of products $\prod_{i=1}^n (A - \lambda I_n)_{i,\sigma(i)}$.
The only permutation contributing a term of degree~$n$ is the identity
$\sigma = \Id$, which gives $\prod_{i=1}^n (a_{ii} - \lambda)$.  The
leading term is $(-\lambda)^n = (-1)^n\lambda^n$.

(ii) In the expansion of $\prod_{i=1}^n(a_{ii} - \lambda)$, the coefficient
of $\lambda^{n-1}$ is $(-1)^{n-1}\sum_{i=1}^n a_{ii} = (-1)^{n-1}\tr(A)$.
No other permutation contributes a term of degree~$n-1$ (such a
permutation would need at least two non-diagonal factors, each of degree~$0$
in~$\lambda$, leaving degree at most $n-2$).

(iii) $\chi_A(0) = \det(A - 0\cdot I_n) = \det(A)$.

(iv) $\lambda$ is an eigenvalue if and only if $\det(A - \lambda I_n) = 0$,
i.e.\ $\chi_A(\lambda) = 0$.

(v) $\chi_B(\lambda) = \det(\inv{P}AP - \lambda I_n)
= \det(\inv{P}(A - \lambda I_n)P)
= \det(\inv{P})\det(A-\lambda I_n)\det(P)
= \chi_A(\lambda)$.
\end{proof}

\begin{definition}{Characteristic polynomial of an endomorphism}%
{def:char-poly-endo}
\index{characteristic polynomial!of an endomorphism}
Let $f\in\End(E)$ and let $A = \Mat_\cB(f)$ be the matrix of~$f$ in
some basis~$\cB$ of~$E$.  The \emph{characteristic polynomial}
of~$f$ is
\[
  \chi_f(\lambda) \deff \chi_A(\lambda) = \det(A - \lambda I_n).
\]
By property~(v) above, this is independent of the choice of
basis~$\cB$.
\end{definition}

\begin{example}{Characteristic polynomial of a $2\times 2$ matrix}%
{ex:char-poly-2x2}
Let $A = \begin{pmatrix} a & b \\ c & d \end{pmatrix}$.  Then
\[
  \chi_A(\lambda) = \det\begin{pmatrix} a-\lambda & b \\ c & d-\lambda
  \end{pmatrix}
  = (a-\lambda)(d-\lambda) - bc
  = \lambda^2 - (a+d)\lambda + (ad-bc)
  = \lambda^2 - \tr(A)\lambda + \det(A).
\]
\end{example}

\begin{example}{Characteristic polynomial of a $3\times 3$ matrix}%
{ex:char-poly-3x3}
Let $A = \begin{pmatrix} 2 & 1 & 0 \\ 0 & 3 & 1 \\ 0 & 0 & 2
\end{pmatrix}$.  Then
\begin{align*}
  \chi_A(\lambda) &= \det\begin{pmatrix} 2-\lambda & 1 & 0 \\
  0 & 3-\lambda & 1 \\ 0 & 0 & 2-\lambda \end{pmatrix} \\
  &= (2-\lambda)(3-\lambda)(2-\lambda) \\
  &= (2-\lambda)^2(3-\lambda)
  = -(\lambda-2)^2(\lambda-3).
\end{align*}
The eigenvalues are $\lambda_1 = 2$ (a double root) and $\lambda_2 = 3$
(a simple root).
\end{example}


% ============================================================
\section{The Cayley--Hamilton theorem}\label{sec:cayley-hamilton}
% ============================================================

\begin{theorem}{Cayley--Hamilton}{thm:cayley-hamilton}
\index{Cayley--Hamilton theorem}
Let $A\in\Mat_n(\K)$ with characteristic polynomial~$\chi_A$.  Then
\[
  \chi_A(A) = 0_{n\times n}.
\]
That is, every square matrix satisfies its own characteristic
polynomial.  Equivalently, if $f\in\End(E)$, then $\chi_f(f) = 0$
(the zero endomorphism).
\end{theorem}

\begin{proof}
Write $B(\lambda) \deff \operatorname{adj}(A - \lambda I_n)$ for the
classical adjoint (adjugate) of $A - \lambda I_n$.  By the adjugate
formula for determinants,
\[
  (A - \lambda I_n)\, B(\lambda) = \det(A - \lambda I_n)\, I_n
  = \chi_A(\lambda)\, I_n.
\]
Each entry of $B(\lambda)$ is a cofactor of $A - \lambda I_n$, hence a
polynomial in~$\lambda$ of degree at most $n-1$.  We can therefore
write
\[
  B(\lambda) = B_0 + B_1\lambda + B_2\lambda^2 + \cdots
    + B_{n-1}\lambda^{n-1},
\]
where $B_0,B_1,\ldots,B_{n-1}\in\Mat_n(\K)$ are constant matrices.
Similarly, write the characteristic polynomial as
\[
  \chi_A(\lambda) = c_0 + c_1\lambda + c_2\lambda^2 + \cdots
    + c_n\lambda^n,
\]
where $c_n = (-1)^n$.

Now expand the identity $(A - \lambda I_n)\,B(\lambda) = \chi_A(\lambda)\,I_n$:
\begin{align*}
  A B_0 &= c_0 I_n, \\
  A B_1 - B_0 &= c_1 I_n, \\
  A B_2 - B_1 &= c_2 I_n, \\
  &\;\;\vdots \\
  A B_{n-1} - B_{n-2} &= c_{n-1} I_n, \\
  -B_{n-1} &= c_n I_n.
\end{align*}
Multiply the $k$-th equation (from $k=0$ to $k=n$) on the left by
$A^k$ and sum:
\begin{align*}
  &\phantom{{}={}} A^0(AB_0) + A^1(AB_1 - B_0) + A^2(AB_2 - B_1)
    + \cdots + A^n(-B_{n-1}) \\
  &= c_0 I_n + c_1 A + c_2 A^2 + \cdots + c_n A^n \\
  &= \chi_A(A).
\end{align*}
On the left-hand side, all terms cancel telescopically:
\[
  AB_0 + A^2 B_1 - AB_0 + A^3 B_2 - A^2 B_1 + \cdots - A^n B_{n-1}
  = 0_{n\times n}.
\]
Therefore $\chi_A(A) = 0_{n\times n}$.
\end{proof}

\begin{remark}{Warning about ``substituting $A$ for $\lambda$''}%
{rem:cayley-hamilton-warning}
One cannot prove Cayley--Hamilton by simply writing
``$\chi_A(A) = \det(A - A\cdot I_n) = \det(0) = 0$.''
This is \emph{fallacious}: $\chi_A(\lambda)$ is a polynomial
in the \emph{scalar}~$\lambda$, while the substitution $\lambda = A$
places a \emph{matrix} inside a determinant in an invalid way.  The
correct proof, as given above, requires careful bookkeeping with matrix
polynomials.
\end{remark}

\begin{example}{Cayley--Hamilton for a $2\times 2$ matrix}{ex:cayley-hamilton-2x2}
Take $A = \begin{pmatrix} 1 & 2 \\ 3 & 4 \end{pmatrix}$.  Then
$\chi_A(\lambda) = \lambda^2 - 5\lambda - 2$.  We verify:
\[
  A^2 - 5A - 2I_2
  = \begin{pmatrix}7 & 10 \\ 15 & 22\end{pmatrix}
  - \begin{pmatrix}5 & 10 \\ 15 & 20\end{pmatrix}
  - \begin{pmatrix}2 & 0 \\ 0 & 2\end{pmatrix}
  = \begin{pmatrix}0 & 0 \\ 0 & 0\end{pmatrix}.\qedhere
\]
\end{example}


% ============================================================
\section{Spectrum of an endomorphism}\label{sec:spectrum}
% ============================================================

\begin{definition}{Spectrum}{def:spectrum}
\index{spectrum}
The \emph{spectrum} of an endomorphism $f\in\End(E)$ (or of a matrix
$A\in\Mat_n(\K)$) is the set of all its eigenvalues:
\[
  \Spec(f) \deff \setcond{\lambda\in\K}{\lambda \text{ is an eigenvalue of } f}.
\]
\end{definition}

\begin{remark}{Dependence on the base field}{rem:spectrum-field}
The spectrum depends on the base field.  For example, the rotation
matrix $R_{\pi/2} = \begin{psmallmatrix} 0 & -1 \\ 1 & 0
\end{psmallmatrix}$ has $\chi_{R_{\pi/2}}(\lambda) = \lambda^2 + 1$,
which has no real roots, so $\Spec_\R(R_{\pi/2}) = \emptyset$.
Over~$\C$, $\Spec_\C(R_{\pi/2}) = \set{i,-i}$.
\end{remark}


% ============================================================
\section{Algebraic and geometric multiplicity}\label{sec:multiplicities}
% ============================================================

\begin{definition}{Algebraic multiplicity}{def:alg-mult}
\index{multiplicity!algebraic}
Let $\lambda$ be an eigenvalue of~$f$.  The \emph{algebraic
multiplicity} of~$\lambda$, denoted $m_a(\lambda)$, is the
multiplicity of~$\lambda$ as a root of the characteristic
polynomial~$\chi_f$.  In other words, $m_a(\lambda)$ is the largest
integer $k$ such that $(\lambda_0 - \lambda)^k$ divides
$\chi_f(\lambda_0)$, where we view $\chi_f$ as a polynomial in the
variable~$\lambda_0$.
\end{definition}

\begin{definition}{Geometric multiplicity}{def:geom-mult}
\index{multiplicity!geometric}
The \emph{geometric multiplicity} of an eigenvalue~$\lambda$,
denoted $m_g(\lambda)$, is the dimension of the eigenspace:
\[
  m_g(\lambda) \deff \dim E_\lambda(f) = \dim\Ker(f - \lambda\Id).
\]
\end{definition}

\begin{theorem}{Inequality between multiplicities}%
{thm:mult-inequality}
\index{multiplicity!inequality}
For every eigenvalue~$\lambda$ of~$f$,
\[
  1 \leq m_g(\lambda) \leq m_a(\lambda).
\]
\end{theorem}

\begin{proof}
The inequality $m_g(\lambda) \geq 1$ holds because $E_\lambda(f)$
contains at least one nonzero eigenvector.

For the inequality $m_g(\lambda) \leq m_a(\lambda)$, let
$d = m_g(\lambda) = \dim E_\lambda(f)$ and choose a basis
$(v_1,\ldots,v_d)$ of~$E_\lambda(f)$.  Extend it to a basis
$\cB = (v_1,\ldots,v_d,v_{d+1},\ldots,v_n)$ of~$E$.  In this basis,
the matrix of~$f$ has the block form
\[
  \Mat_\cB(f) = \begin{pmatrix} \lambda I_d & C \\ 0 & D \end{pmatrix},
\]
where $C\in\Mat_{d,n-d}(\K)$ and $D\in\Mat_{n-d}(\K)$.  This is
because $f(v_i) = \lambda v_i$ for $i = 1,\ldots,d$.  Then
\[
  \chi_f(\mu) = \det\begin{pmatrix} \lambda I_d - \mu I_d & C \\
    0 & D - \mu I_{n-d} \end{pmatrix}
  = (\lambda - \mu)^d \det(D - \mu I_{n-d}).
\]
Therefore $(\lambda - \mu)^d$ divides $\chi_f(\mu)$, which means
$m_a(\lambda) \geq d = m_g(\lambda)$.
\end{proof}

\begin{example}{Strict inequality between multiplicities}%
{ex:strict-mult-inequality}
The matrix $A = \begin{pmatrix} 2 & 1 \\ 0 & 2 \end{pmatrix}$ has
$\chi_A(\lambda) = (2-\lambda)^2$, so $\lambda = 2$ has algebraic
multiplicity $m_a(2) = 2$.  However,
$E_2(A) = \Ker\begin{psmallmatrix} 0 & 1 \\ 0 & 0 \end{psmallmatrix}
= \Span\!\left(\begin{psmallmatrix} 1 \\ 0 \end{psmallmatrix}\right)$,
which has dimension~$1$.  So $m_g(2) = 1 < 2 = m_a(2)$.
\end{example}


% ============================================================
\section{Diagonalizability}\label{sec:diagonalizability}
% ============================================================

\begin{definition}{Diagonalizable endomorphism}{def:diagonalizable-endo}
\index{diagonalizable!endomorphism}
An endomorphism $f\in\End(E)$ is \emph{diagonalizable} if there
exists a basis of~$E$ consisting entirely of eigenvectors of~$f$.
Equivalently, $f$ is diagonalizable if its matrix in some basis is
diagonal.
\end{definition}

\begin{definition}{Diagonalizable matrix}{def:diagonalizable-matrix}
\index{diagonalizable!matrix}
A matrix $A\in\Mat_n(\K)$ is \emph{diagonalizable} if it is similar
to a diagonal matrix, i.e.\ if there exists an invertible matrix
$P\in\GL_n(\K)$ and a diagonal matrix
$D = \diag(\lambda_1,\ldots,\lambda_n)$ such that
\[
  A = P D \inv{P}.
\]
The columns of~$P$ are eigenvectors of~$A$, and the diagonal entries
of~$D$ are the corresponding eigenvalues.
\end{definition}

\begin{center}
\begin{tikzpicture}[>=Stealth, node distance=3.5cm]
  \node (E1) {$E$};
  \node[right of=E1, node distance=4cm] (E2) {$E$};
  \node[below of=E1, node distance=2.5cm] (K1) {$\K^n$};
  \node[below of=E2, node distance=2.5cm] (K2) {$\K^n$};

  \draw[->] (E1) -- node[above] {$f$} (E2);
  \draw[->] (K1) -- node[below] {$D = \diag(\lambda_1,\ldots,\lambda_n)$} (K2);
  \draw[->] (E1) -- node[left] {$[\cdot]_\cB$} (K1);
  \draw[->] (E2) -- node[right] {$[\cdot]_\cB$} (K2);

  \node[below of=K1, node distance=1.5cm, text width=8cm, align=center]
    {\small In the eigenbasis $\cB$, the map $f$ acts as
     multiplication by $\lambda_i$ on each coordinate axis.};
\end{tikzpicture}
\captionof{figure}{Diagonalization: in the eigenbasis $\cB$, the
  endomorphism~$f$ is represented by the diagonal matrix~$D$.}
\label{fig:diag-diagram}
\end{center}

\begin{theorem}{Diagonalization criterion}{thm:diag-criterion}
\index{diagonalizable!criterion}
Let $f\in\End(E)$ with $\dim E = n$, and let
$\Spec(f) = \set{\lambda_1,\ldots,\lambda_r}$.  The following
conditions are equivalent:
\begin{enumerate}[label=(\roman*)]
  \item $f$ is diagonalizable.
  \item $E = E_{\lambda_1}(f) \oplus E_{\lambda_2}(f) \oplus \cdots
    \oplus E_{\lambda_r}(f)$.
  \item $\sum_{i=1}^{r} \dim E_{\lambda_i}(f) = n$.
  \item The characteristic polynomial $\chi_f$ splits over~$\K$
    (i.e.\ has all its roots in~$\K$), \emph{and}
    $m_g(\lambda_i) = m_a(\lambda_i)$ for every eigenvalue
    $\lambda_i$.
\end{enumerate}
\end{theorem}

\begin{proof}
\textbf{(i) $\Rightarrow$ (ii).}\;
If $f$ is diagonalizable, there is a basis $(v_1,\ldots,v_n)$ of~$E$
with $f(v_j) = \mu_j v_j$ for some $\mu_j\in\K$.  Each $v_j$ belongs
to $E_{\mu_j}(f)$.  Since the $v_j$ Span~$E$, we have
$E = \sum_{i=1}^r E_{\lambda_i}(f)$.  This sum is direct by

ef{prop:prop:eigenspace-direct-sum}.

\textbf{(ii) $\Rightarrow$ (iii).}\;
If the direct sum decomposition holds, then
$n = \dim E = \sum_{i=1}^r \dim E_{\lambda_i}(f)$ by the dimension
formula for direct sums.

\textbf{(iii) $\Rightarrow$ (iv).}\;
From $\sum m_g(\lambda_i) = n$ and the inequality
$m_g(\lambda_i) \leq m_a(\lambda_i)$ for each~$i$, together with
$\sum m_a(\lambda_i) \leq n$ (the sum of algebraic multiplicities is
at most~$\deg\chi_f = n$), we get
\[
  n = \sum_{i=1}^r m_g(\lambda_i) \leq \sum_{i=1}^r m_a(\lambda_i) \leq n.
\]
All inequalities are equalities.  In particular,
$m_g(\lambda_i) = m_a(\lambda_i)$ for each~$i$, and
$\sum m_a(\lambda_i) = n$, which means $\chi_f$ splits over~$\K$
(every root is in~$\K$ and the multiplicities account for the full
degree~$n$).

\textbf{(iv) $\Rightarrow$ (i).}\;
If $\chi_f$ splits and $m_g(\lambda_i) = m_a(\lambda_i)$ for each~$i$,
then
\[
  \sum_{i=1}^r \dim E_{\lambda_i}(f)
  = \sum_{i=1}^r m_g(\lambda_i)
  = \sum_{i=1}^r m_a(\lambda_i) = n.
\]
Since the sum of eigenspaces is direct (
ef{prop:prop:eigenspace-direct-sum}),
we have $E = \bigoplus_{i=1}^r E_{\lambda_i}(f)$.  Choosing a basis
for each $E_{\lambda_i}(f)$ and concatenating gives a basis of~$E$
consisting of eigenvectors, so $f$ is diagonalizable.
\end{proof}

\begin{corollary}{Endomorphism with $n$ distinct eigenvalues}%
{cor:n-distinct-eigenvalues}
If $f\in\End(E)$ has $n = \dim E$ distinct eigenvalues, then $f$ is
diagonalizable.
\end{corollary}

\begin{proof}
Each eigenvalue has algebraic multiplicity~$1$ (since there are $n$
eigenvalues and $\deg\chi_f = n$), and the geometric multiplicity
satisfies $1 \leq m_g(\lambda) \leq m_a(\lambda) = 1$.  Thus
$m_g = m_a$ for all eigenvalues and $\chi_f$ splits, so $f$ is
diagonalizable by 
ef{thm:thm:diag-criterion}.
\end{proof}

\begin{remark}{The converse is false}{rem:diag-distinct-converse}
Having $n$ distinct eigenvalues is sufficient but not necessary for
diagonalizability.  For instance, the identity matrix $I_n$ is
diagonal (hence diagonalizable) but has a single eigenvalue
$\lambda = 1$ with multiplicity~$n$.
\end{remark}


% ============================================================
\section{The diagonalization algorithm}\label{sec:diag-algorithm}
% ============================================================

Given $A\in\Mat_n(\K)$, the following procedure determines whether
$A$ is diagonalizable and, if so, produces $P$ and $D$ such that
$A = P D \inv{P}$.

\begin{enumerate}[label=\textbf{Step \arabic*.}]
  \item \textbf{Compute the characteristic polynomial}
    $\chi_A(\lambda) = \det(A - \lambda I_n)$.
  \item \textbf{Find the eigenvalues} by solving $\chi_A(\lambda) = 0$.
    If $\chi_A$ does not split over~$\K$, then $A$ is \emph{not}
    diagonalizable over~$\K$.
  \item \textbf{For each eigenvalue~$\lambda_i$}, compute the eigenspace
    $E_{\lambda_i} = \Ker(A - \lambda_i I_n)$ by solving the
    homogeneous system $(A - \lambda_i I_n)X = \vzero$.
  \item \textbf{Check the multiplicity condition:}
    $\dim E_{\lambda_i} = m_a(\lambda_i)$ for each~$i$.
    If this fails for any eigenvalue, $A$ is \emph{not} diagonalizable.
  \item \textbf{Form} $P$ by placing the eigenvectors (bases of each
    eigenspace) as columns, and $D = \diag(\lambda_1,\ldots,\lambda_1,
    \lambda_2,\ldots,\lambda_2,\ldots)$ with each $\lambda_i$ repeated
    $m_a(\lambda_i)$ times.
\end{enumerate}

\begin{example}{Full diagonalization of a $3\times 3$ matrix}%
{ex:diag-3x3}
\index{diagonalization!example}
Diagonalize $A = \begin{pmatrix} 4 & 0 & 1 \\ 2 & 3 & 2 \\ 1 & 0 & 4
\end{pmatrix}$.

\textbf{Step 1.}  The characteristic polynomial:
\begin{align*}
  \chi_A(\lambda) &= \det\begin{pmatrix} 4-\lambda & 0 & 1 \\
    2 & 3-\lambda & 2 \\ 1 & 0 & 4-\lambda \end{pmatrix} \\
  &= (4-\lambda)\bigl[(3-\lambda)(4-\lambda) - 0\bigr]
    - 0 + 1\bigl[0 - (3-\lambda)\bigr] \\
  &= (4-\lambda)(3-\lambda)(4-\lambda) - (3-\lambda) \\
  &= (3-\lambda)\bigl[(4-\lambda)^2 - 1\bigr] \\
  &= (3-\lambda)(16 - 8\lambda + \lambda^2 - 1) \\
  &= (3-\lambda)(\lambda^2 - 8\lambda + 15) \\
  &= (3-\lambda)(\lambda - 3)(\lambda - 5) \\
  &= -(3-\lambda)^2(\lambda - 5) = -(\lambda-3)^2(\lambda-5).
\end{align*}

\textbf{Step 2.}  Eigenvalues: $\lambda_1 = 3$ (algebraic
multiplicity~$2$) and $\lambda_2 = 5$ (algebraic multiplicity~$1$).

\textbf{Step 3.}  Eigenspaces:

For $\lambda_1 = 3$: solve $(A - 3I_3)X = \vzero$:
\[
  A - 3I_3 = \begin{pmatrix} 1 & 0 & 1 \\ 2 & 0 & 2 \\ 1 & 0 & 1
  \end{pmatrix} \longrightarrow \begin{pmatrix} 1 & 0 & 1 \\ 0 & 0 & 0
  \\ 0 & 0 & 0 \end{pmatrix}.
\]
So $x_1 = -x_3$, $x_2$ free, $x_3$ free.  A basis of $E_3(A)$ is
\[
  \left\{ \begin{pmatrix} 0\\1\\0 \end{pmatrix},\;
  \begin{pmatrix} -1\\0\\1 \end{pmatrix} \right\},
  \quad \dim E_3(A) = 2 = m_a(3).
\]

For $\lambda_2 = 5$: solve $(A - 5I_3)X = \vzero$:
\[
  A - 5I_3 = \begin{pmatrix} -1 & 0 & 1 \\ 2 & -2 & 2 \\ 1 & 0 & -1
  \end{pmatrix} \longrightarrow \begin{pmatrix} 1 & 0 & -1 \\ 0 & 1 & -2
  \\ 0 & 0 & 0 \end{pmatrix}.
\]
So $x_1 = x_3$, $x_2 = 2x_3$, $x_3$ free.  A basis of $E_5(A)$ is
$\left\{ \begin{psmallmatrix} 1\\2\\1 \end{psmallmatrix} \right\}$,
and $\dim E_5(A) = 1 = m_a(5)$.

\textbf{Step 4.}  Since $m_g = m_a$ for both eigenvalues, $A$ is
diagonalizable.

\textbf{Step 5.}  We obtain
\[
  P = \begin{pmatrix} 0 & -1 & 1 \\ 1 & 0 & 2 \\ 0 & 1 & 1
  \end{pmatrix}, \qquad
  D = \begin{pmatrix} 3 & 0 & 0 \\ 0 & 3 & 0 \\ 0 & 0 & 5
  \end{pmatrix}, \qquad
  A = P D \inv{P}.
\]
\end{example}

\begin{example}{A non-diagonalizable matrix}{ex:non-diag}
\index{diagonalizable!non-example}
Consider $A = \begin{pmatrix} 2 & 1 \\ 0 & 2 \end{pmatrix}$.  Then
$\chi_A(\lambda) = (2-\lambda)^2$, so $\lambda = 2$ is the only
eigenvalue with $m_a(2) = 2$.  But
$E_2(A) = \Ker\begin{psmallmatrix} 0 & 1 \\ 0 & 0 \end{psmallmatrix}
= \Span\!\left(\begin{psmallmatrix} 1 \\ 0 \end{psmallmatrix}\right)$,
so $m_g(2) = 1 < 2 = m_a(2)$.  By

ef{thm:thm:diag-criterion}, $A$ is \emph{not} diagonalizable.
\end{example}


% ============================================================
\section{Trigonalization}\label{sec:trigonalization}
% ============================================================

\begin{definition}{Trigonalizable endomorphism}{def:trigonalizable}
\index{trigonalizable}\index{triangularizable|see{trigonalizable}}
An endomorphism $f\in\End(E)$ is \emph{trigonalizable} (over~$\K$)
if there exists a basis of~$E$ in which the matrix of~$f$ is upper
triangular.  A matrix $A\in\Mat_n(\K)$ is trigonalizable if it is
similar to an upper triangular matrix.
\end{definition}

\begin{theorem}{Trigonalization over an algebraically closed field}%
{thm:trigonalization}
\index{trigonalization theorem}
Let $f\in\End(E)$ with $\dim E = n$.  Then $f$ is trigonalizable
over~$\K$ if and only if $\chi_f$ splits over~$\K$.

In particular, every endomorphism of a finite-dimensional complex
vector space is trigonalizable (since $\C$ is algebraically closed).
\end{theorem}

\begin{proof}
$(\Leftarrow)$\; We proceed by induction on $n = \dim E$.
The case $n = 1$ is trivial.

Since $\chi_f$ splits, $f$ has at least one eigenvalue
$\lambda_1\in\K$.  Let $v_1$ be a corresponding eigenvector and set
$F = \Span(v_1)$.  The quotient space $\bar{E} = E/F$ has
dimension~$n-1$, and $f$ induces a well-defined endomorphism
$\bar{f}\colon\bar{E}\to\bar{E}$ (since $F$ is $f$-invariant).
Moreover, $\chi_{\bar{f}}$ divides $\chi_f$ (up to a constant),
so $\chi_{\bar{f}}$ also splits over~$\K$.

By the induction hypothesis, $\bar{f}$ is trigonalizable: there
exists a basis $(\bar{v}_2,\ldots,\bar{v}_n)$ of~$\bar{E}$ in which
$\bar{f}$ is upper triangular.  Choosing representatives
$v_2,\ldots,v_n\in E$ with $\bar{v}_i = v_i + F$, the family
$(v_1,v_2,\ldots,v_n)$ is a basis of~$E$ and the matrix of~$f$ in
this basis is upper triangular with $\lambda_1$ in the $(1,1)$
position.

$(\Rightarrow)$\; If $f$ is trigonalizable with upper triangular
matrix $T$, then $\chi_f(\lambda) = \chi_T(\lambda) = \prod_{i=1}^n
(t_{ii} - \lambda)$, which splits over~$\K$.
\end{proof}

\begin{corollary}{Trigonalization over $\C$}{cor:trig-C}
Every matrix $A\in\Mat_n(\C)$ is trigonalizable over~$\C$.
\end{corollary}


% ============================================================
\section{Minimal polynomial}\label{sec:minimal-poly}
% ============================================================

\begin{definition}{Annihilating polynomial}{def:annihilating-poly}
\index{annihilating polynomial}
A polynomial $p\in\K[\lambda]$ is an \emph{annihilating polynomial}
(or \emph{annihilator}) of $f\in\End(E)$ if $p(f) = 0$ (the zero
endomorphism).  Similarly, $p$ annihilates $A\in\Mat_n(\K)$ if
$p(A) = 0_{n\times n}$.
\end{definition}

By the Cayley--Hamilton theorem (
ef{thm:thm:cayley-hamilton}), $\chi_f$
is always an annihilating polynomial of~$f$.  The minimal polynomial
is the ``smallest'' such polynomial.

\begin{definition}{Minimal polynomial}{def:minimal-poly}
\index{minimal polynomial}
The \emph{minimal polynomial} of $f\in\End(E)$, denoted $\mu_f$, is
the unique monic polynomial of smallest degree that annihilates~$f$.
\end{definition}

\begin{proposition}{Existence and uniqueness of the minimal polynomial}%
{prop:minimal-poly-unique}
The minimal polynomial $\mu_f$ exists and is unique.  Moreover, $\mu_f$
divides every annihilating polynomial of~$f$.
\end{proposition}

\begin{proof}
\textit{Existence.}\; The set of annihilating polynomials is nonempty
(it contains $\chi_f$ by Cayley--Hamilton) and every polynomial has a
nonnegative degree, so there is one of minimal degree.  Normalising
to be monic gives~$\mu_f$.

\textit{Divides every annihilator.}\; Let $p$ be an annihilating
polynomial.  Perform Euclidean division: $p = q\mu_f + r$ with
$\deg r < \deg\mu_f$.  Then $r(f) = p(f) - q(f)\mu_f(f) = 0 - 0 = 0$.
If $r \neq 0$, then $r/\operatorname{lc}(r)$ would be a monic
annihilator of smaller degree than~$\mu_f$, a contradiction.  Hence
$r = 0$ and $\mu_f \mid p$.

\textit{Uniqueness.}\; If $\mu$ and $\mu'$ are both monic annihilating
polynomials of minimal degree, then $\mu \mid \mu'$ and $\mu' \mid \mu$
(by the above), so $\mu = \mu'$ (both being monic).
\end{proof}

\begin{proposition}{Properties of the minimal polynomial}%
{prop:minimal-poly-properties}
Let $f\in\End(E)$.  Then:
\begin{enumerate}[label=(\roman*)]
  \item $\mu_f$ divides $\chi_f$.
  \item $\mu_f$ and $\chi_f$ have the same roots (i.e.\ the same
    eigenvalues), though possibly with different multiplicities.
  \item If $A$ and $B$ are similar matrices, they have the same
    minimal polynomial.
\end{enumerate}
\end{proposition}

\begin{proof}
(i) Cayley--Hamilton gives $\chi_f(f) = 0$, so $\mu_f \mid \chi_f$.

(ii) Since $\mu_f \mid \chi_f$, every root of $\mu_f$ is a root of
$\chi_f$.  Conversely, if $\lambda$ is an eigenvalue of~$f$ with
eigenvector~$v$, then $\mu_f(f)(v) = \mu_f(\lambda)v = \vzero$,
so $\mu_f(\lambda) = 0$ (since $v \neq \vzero$).

(iii) If $B = \inv{P}AP$, then $p(B) = \inv{P}p(A)P$ for any
polynomial~$p$.  Thus $p(A) = 0$ if and only if $p(B) = 0$.
\end{proof}

\begin{theorem}{Diagonalizability via the minimal polynomial}%
{thm:diag-minimal-poly}
\index{diagonalizable!minimal polynomial criterion}
An endomorphism $f\in\End(E)$ is diagonalizable if and only if its
minimal polynomial $\mu_f$ splits into distinct linear factors
over~$\K$:
\[
  \mu_f(\lambda) = (\lambda - \lambda_1)(\lambda - \lambda_2)\cdots
    (\lambda - \lambda_r)
\]
where $\lambda_1,\ldots,\lambda_r$ are pairwise distinct.
\end{theorem}

\begin{proof}
$(\Rightarrow)$\; Suppose $f$ is diagonalizable with eigenvalues
$\lambda_1,\ldots,\lambda_r$.  Set
$p(\lambda) = (\lambda-\lambda_1)\cdots(\lambda-\lambda_r)$.  Since $f$
is diagonalizable, there is a basis of eigenvectors.  For any
eigenvector~$v$ with eigenvalue~$\lambda_i$,
$p(f)(v) = \prod_{j\neq i}(\lambda_i - \lambda_j)\cdot 0 \cdot v = \vzero$.
Wait---more carefully:
$p(f)(v) = (f-\lambda_1\Id)\cdots(f-\lambda_r\Id)(v)$.  Since
$(f - \lambda_i\Id)(v) = \vzero$, the entire product gives~$\vzero$.
As this holds for every eigenvector and eigenvectors Span~$E$, we get
$p(f) = 0$.  So $\mu_f \mid p$.  Since the roots of $\mu_f$ include
all eigenvalues (by 
ef{prop:prop:minimal-poly-properties}(ii)) and
$\mu_f$ divides a product of distinct linear factors, $\mu_f$ itself
is a product of distinct linear factors.

$(\Leftarrow)$\; Suppose $\mu_f(\lambda)=(\lambda-\lambda_1)\cdots
(\lambda-\lambda_r)$ with distinct $\lambda_i$.  We prove
$E = \bigoplus_{i=1}^r E_{\lambda_i}(f)$ by induction on~$r$.

For $r = 1$: $\mu_f(\lambda) = \lambda - \lambda_1$, so
$f = \lambda_1\Id$ and every nonzero vector is an eigenvector.
Thus $E = E_{\lambda_1}(f)$.

For the inductive step, write
$\mu_f = (\lambda - \lambda_r)\,q(\lambda)$ where
$q(\lambda) = \prod_{i=1}^{r-1}(\lambda - \lambda_i)$.  Since
$\gcd(q, \lambda-\lambda_r) = 1$ (the $\lambda_i$ are distinct),
there exist polynomials $a,b\in\K[\lambda]$ with
$a(\lambda)\,q(\lambda) + b(\lambda)(\lambda-\lambda_r) = 1$.
Substituting $f$:
\[
  a(f)\circ q(f) + b(f)\circ(f - \lambda_r\Id) = \Id.
\]
For any $v\in E$, write $v = a(f)(q(f)(v)) + b(f)((f-\lambda_r\Id)(v))$.
Let $u = q(f)(v)$ and $w = (f-\lambda_r\Id)(v)$.  Then
$(f-\lambda_r\Id)(u) = (f-\lambda_r\Id)(q(f)(v)) = \mu_f(f)(v) = \vzero$,
so $u\in E_{\lambda_r}(f)$.  Similarly,
$q(f)(w) = q(f)(f-\lambda_r\Id)(v) = \mu_f(f)(v) = \vzero$, so
$w\in\Ker(q(f))$.

Now $q(f)$ acts on $F = \Ker(q(f))$, and $q$ annihilates $f\restrict{F}$.
So $\mu_{f|_F}$ divides $q = \prod_{i=1}^{r-1}(\lambda-\lambda_i)$.
By the induction hypothesis, $F = \bigoplus_{i=1}^{r-1}E_{\lambda_i}(f|_F)
\subset \bigoplus_{i=1}^{r-1}E_{\lambda_i}(f)$.

Since $v = a(f)(u) + b(f)(w)$ and $u\in E_{\lambda_r}(f)$,
$w\in F\subset\bigoplus_{i=1}^{r-1}E_{\lambda_i}(f)$, we see that
every $v\in E$ lies in $\sum_{i=1}^r E_{\lambda_i}(f)$.  The sum is
direct by 
ef{prop:prop:eigenspace-direct-sum}.
\end{proof}


% ============================================================
\section{Applications}\label{sec:eigen-applications}
% ============================================================

% -------------------------
\subsection{Computing powers of a matrix}\label{subsec:matrix-powers}
% -------------------------

\begin{proposition}{Powers of a diagonalizable matrix}%
{prop:matrix-power}
\index{matrix power}
If $A = P D \inv{P}$ where $D = \diag(\lambda_1,\ldots,\lambda_n)$,
then for every integer $k\geq 0$,
\[
  A^k = P\,\diag(\lambda_1^k,\ldots,\lambda_n^k)\,\inv{P}.
\]
\end{proposition}

\begin{proof}
By induction, $A^k = (PD\inv{P})^k = PD^k\inv{P}$, since
$(PD\inv{P})(PD\inv{P}) = PD(\inv{P}P)D\inv{P} = PD^2\inv{P}$,
and so on.  Moreover, $D^k = \diag(\lambda_1^k,\ldots,\lambda_n^k)$.
\end{proof}

% -------------------------
\subsection{Linear recurrences: the Fibonacci sequence}%
\label{subsec:fibonacci}
% -------------------------

\begin{example}{Fibonacci sequence via diagonalization}{ex:fibonacci}
\index{Fibonacci sequence}\index{linear recurrence}
The Fibonacci sequence is defined by $F_0 = 0$, $F_1 = 1$, and
$F_{n+2} = F_{n+1} + F_n$.  Set
$V_n = \begin{psmallmatrix} F_{n+1} \\ F_n \end{psmallmatrix}$.  Then
$V_{n+1} = A V_n$ where
$A = \begin{psmallmatrix} 1 & 1 \\ 1 & 0 \end{psmallmatrix}$.

The characteristic polynomial is $\chi_A(\lambda) = \lambda^2 - \lambda - 1$,
with roots $\vphi = \frac{1+\sqrt{5}}{2}$ and
$\psi = \frac{1-\sqrt{5}}{2}$.

Eigenvectors:
$v_1 = \begin{psmallmatrix}\vphi\\1\end{psmallmatrix}$ for $\vphi$
and $v_2 = \begin{psmallmatrix}\psi\\1\end{psmallmatrix}$ for $\psi$.
So $P = \begin{psmallmatrix}\vphi & \psi \\ 1 & 1\end{psmallmatrix}$
and
$A = P\begin{psmallmatrix}\vphi & 0 \\ 0 & \psi\end{psmallmatrix}\inv{P}$.

From $V_n = A^n V_0$ and $V_0 = \begin{psmallmatrix}1\\0\end{psmallmatrix}$,
we extract:
\[
  F_n = \frac{\vphi^n - \psi^n}{\vphi - \psi}
  = \frac{1}{\sqrt{5}}\left[\left(\frac{1+\sqrt{5}}{2}\right)^{\!n}
    - \left(\frac{1-\sqrt{5}}{2}\right)^{\!n}\right].
\]
This is \emph{Binet's formula}\index{Binet's formula}.
\end{example}

% -------------------------
\subsection{Markov chains and steady-state distributions}%
\label{subsec:markov}
% -------------------------

\begin{example}{Steady-state of a Markov chain}{ex:markov}
\index{Markov chain}\index{steady-state distribution}
A weather model has two states: Sunny (S) and Rainy (R), with
transition matrix
\[
  T = \begin{pmatrix} 0.8 & 0.4 \\ 0.2 & 0.6 \end{pmatrix},
\]
where column~$j$ gives the transition probabilities from state~$j$.
If $\vec{p}_n$ is the state vector at time~$n$, then
$\vec{p}_{n+1} = T\vec{p}_n$, so $\vec{p}_n = T^n\vec{p}_0$.

The eigenvalues of~$T$ are $\lambda_1 = 1$ and $\lambda_2 = 0.4$.
The eigenvector for $\lambda_1 = 1$ (normalised to sum to~$1$) is
\[
  \vec{\pi} = \begin{pmatrix} 2/3 \\ 1/3 \end{pmatrix}.
\]
Since $|\lambda_2| < 1$, as $n\to\infty$ the component along the
second eigenvector decays, and
$\vec{p}_n \to \vec{\pi}$ regardless of~$\vec{p}_0$.  Thus the
long-run probability of sunny weather is~$2/3$.
\end{example}

% -------------------------
\subsection{Systems of linear ODEs}\label{subsec:odes}
% -------------------------

\begin{example}{System of ODEs via diagonalization}{ex:odes}
\index{differential equation!system}
Consider the system $\vec{x}'(t) = A\vec{x}(t)$ with
$A = \begin{psmallmatrix} -3 & 1 \\ 1 & -3 \end{psmallmatrix}$.
The eigenvalues are $\lambda_1 = -2$ (eigenvector
$v_1 = \begin{psmallmatrix}1\\1\end{psmallmatrix}$) and
$\lambda_2 = -4$ (eigenvector
$v_2 = \begin{psmallmatrix}1\\-1\end{psmallmatrix}$).

Setting $\vec{x}(t) = P\vec{y}(t)$ where
$P = \begin{psmallmatrix}1 & 1 \\ 1 & -1\end{psmallmatrix}$,
the system decouples to $\vec{y}'(t) = D\vec{y}(t)$ with
$D = \diag(-2,-4)$.  The solution is
\[
  \vec{y}(t) = \begin{pmatrix} c_1 e^{-2t} \\ c_2 e^{-4t} \end{pmatrix},
  \qquad
  \vec{x}(t) = c_1 e^{-2t}\begin{pmatrix}1\\1\end{pmatrix}
    + c_2 e^{-4t}\begin{pmatrix}1\\-1\end{pmatrix},
\]
where $c_1,c_2$ are determined by the initial condition
$\vec{x}(0)$.
\end{example}


% ============================================================
\section{Exercises}\label{sec:eigenvalue-exercises}
% ============================================================

\begin{exercise}{Eigenvalues of a $2\times 2$ matrix \basic}%
{exo:eigen-2x2-basic}
Find the eigenvalues and eigenvectors of
$A = \begin{pmatrix} 5 & 4 \\ 1 & 2 \end{pmatrix}$.
Is $A$ diagonalizable?  If so, find $P$ and $D$ such that
$A = PD\inv{P}$.
\end{exercise}

\begin{exercise}{Eigenvalues of a triangular matrix \basic}%
{exo:triangular-eigen}
Let $A\in\Mat_n(\K)$ be upper triangular.  Show that the eigenvalues
of~$A$ are exactly the diagonal entries $a_{11},a_{22},\ldots,a_{nn}$.
\end{exercise}

\begin{exercise}{Eigenvalues of a projection \basic}{exo:proj-eigen}
Let $p\in\End(E)$ be a projection, i.e.\ $p^2 = p$.  Show that the
only possible eigenvalues of~$p$ are $0$ and~$1$, and that $p$ is
diagonalizable.

\noindent\textit{Hint:} What is the minimal polynomial of~$p$?
\end{exercise}

\begin{exercise}{Eigenvalues and trace/determinant \basic}%
{exo:trace-det-eigen}
Let $A\in\Mat_2(\K)$ with eigenvalues $\lambda_1,\lambda_2$.  Show
that $\tr(A) = \lambda_1 + \lambda_2$ and
$\det(A) = \lambda_1\lambda_2$.  Generalise to $n\times n$ matrices.
\end{exercise}

\begin{exercise}{Eigenvalues of $A^2$, $A^{-1}$, $A + cI$ \intermediate}%
{exo:eigen-transformations}
Let $A\in\Mat_n(\K)$ with eigenvalue $\lambda$ and eigenvector~$v$.
\begin{enumerate}[label=(\alph*)]
  \item Show that $\lambda^2$ is an eigenvalue of $A^2$ with the same
    eigenvector~$v$.  Generalise to $A^k$.
  \item If $A$ is invertible, show that $\lambda^{-1}$ is an eigenvalue
    of~$\inv{A}$.
  \item Show that $\lambda + c$ is an eigenvalue of $A + cI$ for
    any $c\in\K$.
\end{enumerate}
\end{exercise}

\begin{exercise}{Diagonalization of a $3\times 3$ matrix \intermediate}%
{exo:diag-3x3}
Diagonalize the matrix
$A = \begin{pmatrix} 1 & 2 & 0 \\ 0 & 3 & 0 \\ 2 & 1 & -1 \end{pmatrix}$.
\end{exercise}

\begin{exercise}{Non-diagonalizable matrix \intermediate}{exo:non-diag}
Show that $A = \begin{pmatrix} 1 & 1 & 0 \\ 0 & 1 & 1 \\
0 & 0 & 1 \end{pmatrix}$ is not diagonalizable over any field.
What is its minimal polynomial?
\end{exercise}

\begin{exercise}{Cayley--Hamilton application \intermediate}%
{exo:cayley-hamilton-app}
Let $A = \begin{pmatrix} 1 & 2 \\ 3 & 4 \end{pmatrix}$.
\begin{enumerate}[label=(\alph*)]
  \item Find the characteristic polynomial $\chi_A$.
  \item Use the Cayley--Hamilton theorem to express $\inv{A}$ as a
    polynomial in~$A$ (i.e.\ in the form $aI + bA$).
  \item Compute $A^5$ using Cayley--Hamilton to reduce to lower
    powers.
\end{enumerate}
\end{exercise}

\begin{exercise}{Powers via diagonalization \intermediate}{exo:powers}
\index{matrix power!via diagonalization}
Let $A = \begin{pmatrix} 2 & 1 \\ 0 & 3 \end{pmatrix}$.
\begin{enumerate}[label=(\alph*)]
  \item Diagonalize $A$.
  \item Compute $A^{100}$.
\end{enumerate}
\end{exercise}

\begin{exercise}{Simultaneous diagonalization \challenging}%
{exo:simul-diag}
\index{simultaneous diagonalization}
Let $A,B\in\Mat_n(\K)$ be diagonalizable matrices that commute
($AB = BA$).  Show that $A$ and $B$ are \emph{simultaneously
diagonalizable}: there exists an invertible $P$ such that both
$\inv{P}AP$ and $\inv{P}BP$ are diagonal.

\noindent\textit{Hint:} Show that eigenspaces of~$A$ are invariant
under~$B$, then diagonalize the restriction of~$B$ to each eigenspace.
\end{exercise}

\begin{exercise}{Computing the Fibonacci sequence \intermediate}%
{exo:fibonacci}
\begin{enumerate}[label=(\alph*)]
  \item Diagonalize $A = \begin{psmallmatrix} 1 & 1 \\ 1 & 0
    \end{psmallmatrix}$ over~$\R$.
  \item Derive Binet's formula:
    $F_n = \dfrac{1}{\sqrt{5}}\!\left[\left(\dfrac{1+\sqrt{5}}{2}
    \right)^{\!n} - \left(\dfrac{1-\sqrt{5}}{2}\right)^{\!n}\right]$.
  \item Show that $F_n$ is the nearest integer to
    $\dfrac{1}{\sqrt{5}}\left(\dfrac{1+\sqrt{5}}{2}\right)^{\!n}$
    for all $n\geq 0$.
\end{enumerate}
\end{exercise}

\begin{exercise}{Minimal polynomial \challenging}{exo:minimal-poly}
For each of the following matrices, find the minimal polynomial and
determine whether the matrix is diagonalizable.
\begin{enumerate}[label=(\alph*)]
  \item $A = \begin{pmatrix} 3 & 0 & 0 \\ 0 & 3 & 0 \\ 0 & 0 & 5
    \end{pmatrix}$
  \item $B = \begin{pmatrix} 3 & 1 & 0 \\ 0 & 3 & 0 \\ 0 & 0 & 5
    \end{pmatrix}$
  \item $C = \begin{pmatrix} 3 & 1 & 0 \\ 0 & 3 & 1 \\ 0 & 0 & 3
    \end{pmatrix}$
\end{enumerate}
\end{exercise}

\begin{exercise}{Markov chain steady state \intermediate}{exo:markov}
\index{Markov chain!exercise}
A rat in a maze has three rooms.  At each time step, it moves
according to the transition matrix
\[
  T = \begin{pmatrix} 0.5 & 0.25 & 0.25 \\ 0.25 & 0.5 & 0.25 \\
  0.25 & 0.25 & 0.5 \end{pmatrix}.
\]
\begin{enumerate}[label=(\alph*)]
  \item Find the eigenvalues and eigenspaces of~$T$.
  \item Determine the steady-state distribution $\vec{\pi}$ satisfying
    $T\vec{\pi} = \vec{\pi}$, $\sum_i \pi_i = 1$.
  \item Compute $T^n$ explicitly and verify that
    $T^n \to \vec{\pi}\,\transpose{\vec{1}}$ as $n\to\infty$ (where
    $\vec{1} = (1,1,1)^{\mathsf{T}}/3$).
\end{enumerate}
\end{exercise}


% ============================================================
\section{Chapter summary}\label{sec:eigenvalue-summary}
% ============================================================

\begin{itemize}
  \item An \textbf{eigenvalue} $\lambda$ of $f\in\End(E)$ satisfies
    $f(v) = \lambda v$ for some nonzero~$v$ (the \textbf{eigenvector}).
    The \textbf{eigenspace} $E_\lambda(f) = \Ker(f - \lambda\Id)$ is a
    subspace.

  \item The \textbf{characteristic polynomial}
    $\chi_f(\lambda) = \det(f - \lambda\Id)$ has degree~$n$; its roots
    in~$\K$ are exactly the eigenvalues.  The set of eigenvalues is
    the \textbf{spectrum} $\Spec(f)$.

  \item The \textbf{Cayley--Hamilton theorem}: every matrix satisfies
    its own characteristic polynomial, $\chi_A(A) = 0$.

  \item For each eigenvalue~$\lambda$, the \textbf{geometric
    multiplicity} $m_g(\lambda) = \dim E_\lambda$ satisfies
    $1 \leq m_g(\lambda) \leq m_a(\lambda)$ (the \textbf{algebraic
    multiplicity}).

  \item $f$ is \textbf{diagonalizable} if and only if $\chi_f$ splits
    over~$\K$ and $m_g = m_a$ for every eigenvalue.  Equivalently,
    $E = \bigoplus_{\lambda\in\Spec(f)} E_\lambda(f)$.

  \item \textbf{Diagonalization algorithm}: compute $\chi_f$, find
    eigenvalues, compute eigenspaces, check $m_g = m_a$, form the
    change-of-basis matrix~$P$.

  \item Every endomorphism whose characteristic polynomial splits is
    \textbf{trigonalizable}.  Over~$\C$, every endomorphism is
    trigonalizable.

  \item The \textbf{minimal polynomial} $\mu_f$ is the monic
    polynomial of smallest degree annihilating~$f$.  It divides~$\chi_f$
    and shares its roots.  $f$ is diagonalizable if and only if $\mu_f$
    splits into distinct linear factors.

  \item \textbf{Applications}: if $A = PD\inv{P}$ then
    $A^k = PD^k\inv{P}$, which allows efficient computation of matrix
    powers, closed-form solutions of linear recurrences (Fibonacci),
    steady-state analysis of Markov chains, and explicit solutions of
    linear systems of ODEs.
\end{itemize}
